\chapter{知识的质量评估与修正}
\label{chap:knowledge-quality-assessment-correction}

假设一个数据集含有一万条候选标签，检测程序从中挑出一百条高风险记录，复核后确认其中四十条有错。这个结果说明检测程序找到了值得排查的区域，却不能推出全体标签的错误率是$40\%$：这一百条本来就是按风险选出的，不代表其余九千九百条。即使随后改正了这四十条，也还不能断言修正规则对新数据有效，因为规则可能只是适配了已经看过的问题。质量改进中的三个常见误判由此出现：把重点排查当作总体审计，把检测线索当作错误证据，把已改结果的复核当作修正策略的留出验证。

本章讨论怎样建立评估依据、用有限审计样本估计质量、定位问题、实施修正并验证效果。检查对象既可以是类别、数值和结构化答案，也可以是样本关系、内容恢复目标和训练中生成的表示目标。不同知识具有不同的可核验含义：任务答案可以与独立参考比较，关系知识要检查对象、方向与成立条件，表示目标通常只能先核查构造过程，再用独立下游任务评价用途。把它们都压成一个“正确率”，会掩盖被测对象和证据含义的差异。

为了贯穿各节，沿用第~\ref{chap:knowledge-fusion}章的开发对象$x_2$和知识版本K17。两条同模板规则共享旧词典，都输出“山茶”，模型输出“茶梅”，人工尚未请求；K17采用逐来源等权多数票，因而把同一规则依据重复计票，输出“山茶”。独立照片与采集记录支持“茶梅”，这个对象便成为检查共享来源与融合配置的入口。$x_2$用于讲解和开发诊断，不进入后文的最终留出验收总体。

同一批花卉知识还可能出现局部框漏掉花瓣、正对应关系把两株不同植物配成同一对象，或裁剪恢复目标因关键对象被删除而失去声明语义等问题。审计记录应把知识版本及其来源路径、抽样设计、参考证据和裁决理由连起来。处理阶段取“待复核”或“已完成”，质量结论另记为“正确、错误、部分正确、缺失、不适用、无法裁决”。\term{待复核线索}表示检查触发了怀疑，尚无质量结论；完成审计则既可能确认原记录正确，也可能确认错误，或者因证据不足而无法裁决。“信息不足”“规范歧义”等作为无法裁决的原因保存。

\section{质量评估依据的建立}
\label{sec:knowledge-quality-assessment-correction-basis}

质量测量之前必须先回答“评什么、凭什么判断、结论能推广到哪里”。任务标准规定需要检查的属性，参考依据规定怎样裁决单条记录，抽样方案规定有限证据代表哪一部分总体。这三项共同构成\term{审计方案}（audit plan）。缺少任何一项，计算出的比例都可能有数值却没有明确含义。

\subsection{评估标准的确定}
\label{sec:knowledge-quality-assessment-correction-criteria}

\term{评估标准}把用途要求转成可复核的判断。正确性检查知识内容是否符合任务定义，完整性检查所需对象或局部信息是否缺失，粒度检查答案是否细到任务要求的层次，时效性检查事实和来源版本是否仍适用于目标时间。格式合规只说明记录能被解析，证据充分说明当前材料足以作出判断，语义正确说明结论与参考依据一致，适于下游任务则说明知识在指定使用方式下带来可接受的效果；四者不能互相替代。

评估单位也随知识形式变化。类别答案的单位可以是一条样本标签；对象检测需要把对象存在性、边界和类别分开；关系知识要连同两个对象、关系方向、时间与上下文共同核验；内容预测目标要检查被保留内容是否确实来自同一对象；表示目标要检查视图、教师或聚类目标的生成条件。若任务允许多个合理答案，标准还要说明裁决对象是答案集合、候选排序、偏好关系还是单一类别。

沿用花卉案例，一张照片标为“山茶”可能有四种不同问题：物种类别错误；照片中还有另一株花却未标出；任务要求亚种而答案只到物种；标签引用的是已经废止的分类版本。另一张照片的两个裁剪都来自同一原图，但其中一个裁剪只剩背景，此时配对记录在形式上正确，声明的“保持对象语义”却已经失效。表~\ref{tab:knowledge-quality-assessment-correction-criteria}把这些差异放到同一判断框架中。

\begin{table*}[htbp]
  \centering
  \small
  \begingroup
  \setlength{\tabcolsep}{2.5pt}
  \begin{tabularx}{\linewidth}{@{}>{\RaggedRight\setlength{\parindent}{0pt}\arraybackslash}p{26mm}>{\RaggedRight\setlength{\parindent}{0pt}\arraybackslash}p{42mm}>{\RaggedRight\setlength{\parindent}{0pt}\arraybackslash}p{47mm}>{\RaggedRight\setlength{\parindent}{0pt}\arraybackslash}X@{}}
    \toprule
    \textbf{知识对象} & \textbf{可能问题} & \textbf{独立参考} & \textbf{可测属性} \\
    \midrule
    任务类别 & 错类、强制单解、时效失配 & 原始对象、规范与独立复标 & 参考一致性、错误类型、无法裁决比例 \\
    局部结构 & 漏标、边界偏移、类别错配 & 原图或原文、对象级规范与独立定位 & 对象召回、边界重合与局部完整性 \\
    样本关系 & 错配、漏配、方向或条件错误 & 采集标识、对象身份与时间记录 & 错配率、关系覆盖与条件违例率 \\
    内容预测目标 & 内容泄漏、对象错误、可见范围越界 & 原始观测、遮蔽记录与对象对应 & 目标一致性、泄漏率与对象错配率 \\
    表示目标 & 教师退化、语义被破坏、分组不稳定 & 构造日志、变换条件与冻结的用途评价 & 构造违例、稳定性、覆盖与下游效用 \\
    \bottomrule
  \end{tabularx}
  \endgroup
  \caption{不同知识的质量问题与判断依据。独立参考应与待核验属性相匹配；构造条件检查与下游用途评价分别提供证据，不能用一种结果替代另一种。}
  \label{tab:knowledge-quality-assessment-correction-criteria}
\end{table*}

标准还要约定严重程度、不适用和无法裁决。把物种写错会直接改变监督目标，别名格式不统一可能只影响接口，照片证据不足则不应被迫归入“正确”或“错误”。若规范后来改变了物种边界，旧标签与新标签的差异可能来自评价基准变化，而不是原执行过程退化。数据文档应记录数据组成、采集、处理、用途和维护等背景，正是为了让这些判断条件可追溯\scite{gebru2021datasheets}

\subsection{参考依据的建立}
\label{sec:knowledge-quality-assessment-correction-reference}

\term{参考依据}是裁决当前质量属性所用的独立材料和程序，不等同于再找一个标签覆盖原标签。可以独立核实答案时，参考依据可以来自盲化复标、领域专家裁决、可信原始记录或可执行检查。盲化是指复核者先看原始对象和任务规范，不看原标签、模型建议及其置信分数，以减少被已有答案牵引；独立判断完成后，再展示冲突信息用于定向裁决。

多人意见一致与证据充分也不是同一回事。复核者可能都依赖同一本错误名录，模型和人工也可能共享同一条过时记录。因而参考记录应保存证据来源、适用时间、复核者角色、独立判断、冲突意见和裁决理由。合理多解可以保留答案集合或条件化结论；原始材料不足时应标为无法裁决。为了得到整齐的单标签而强行消除分歧，会把参考集本身变成新的错误来源。

例如，原标签与模型建议都把一张花卉照片判为山茶，但采集编号指向另一株植物。复核者先根据照片和同一植株的其他视图独立判断，再查采集记录，发现照片写回时发生对象错配。这里推翻原结果的不是“第三票”，而是独立对象证据。参考标签仍可能出错，所以高风险结论可以安排二次复核，并保留能够回到原始材料的路径。

样本关系和表示目标常常没有逐条标准答案。对一对视图，可以核对它们是否来自同一对象、变换是否保留任务语义、时间顺序是否满足要求；对教师表示，可以核对教师版本、输入可见范围和更新规则。这些检查能证明声明的构造条件是否执行，却不能证明该关系或表示一定改善学习。用途收益要在冻结的下游任务与评价集上另行验证。

\subsection{审计样本的选择}
\label{sec:knowledge-quality-assessment-correction-sampling}

\term{总体审计}（population audit）用有已知入样机制的样本估计目标总体质量；\term{重点排查}（targeted inspection）则把预算集中在高风险对象上，以便尽快发现和处理问题。两者都需要复核样本，却回答不同问题。前者要求目标总体中的重要区域有被抽中的机会，后者可以故意偏向冲突、高损失或高影响记录，但重点集合中的错误比例只描述该选择条件下的对象。

设目标总体为$U=\{1,\ldots,N\}$，审计样本为$S\subseteq U$，对象$i$的入样概率为$\pi_i=\mathbb{P}(i\in S)$。简单随机无放回抽样令所有$\pi_i=n/N$；分层抽样先按类别、来源、时间或风险区域划分总体，再在各层内随机抽取；按簇抽样则先抽文档、用户或采集批次，并审核簇内记录。不等概率设计可以增加少数类别或高损失区域的份额，但必须保存$\pi_i>0$，后续估计才能恢复目标总体的权重。

图~\ref{fig:knowledge-quality-assessment-correction-audit-sampling}把同一教学总体分别送入分层随机审计和高风险重点排查。读图时应比较两侧被选对象的组成：左侧的层内概率为总体估计保留了设计依据，右侧的选择条件则主动提高了可疑对象占比。

\begin{figure*}[htbp]
  \centering
  % 图7-1：同一教学总体上的概率审计与重点排查，按169 mm全版心宽度设计。
\begin{tikzpicture}[
  x=1mm,
  y=1mm,
  font=\figurefont\footnotesize,
  text=BookInk,
  panel/.style={draw=BookRule,fill=white,line width=0.8pt},
  common item/.style={circle,draw=FigureInput,fill=FigureInput!9,line width=0.7pt,minimum size=4.2mm,inner sep=0pt},
  key item/.style={rectangle,draw=FigureModel,fill=FigureModel!9,line width=0.7pt,minimum size=4.2mm,inner sep=0pt},
  audit ring/.style={circle,draw=BookBlue,line width=1.2pt,minimum size=6.8mm,inner sep=0pt},
  focus ring/.style={circle,draw=BookAmber,line width=1.2pt,dashed,minimum size=6.8mm,inner sep=0pt},
  problem/.style={font=\figurefont\bfseries\scriptsize,text=BookAmber},
  note/.style={font=\figurefont\scriptsize,text=BookMuted,align=center}
]
  \draw[panel] (0,0) rectangle (80,58);
  \draw[panel] (86,0) rectangle (166,58);
  \draw[BookRule,line width=0.6pt,dashed] (2,20) -- (78,20);
  \draw[BookRule,line width=0.6pt,dashed] (88,20) -- (164,20);

  \node[anchor=north west,font=\figurefont\bfseries] at (2,56) {(a) 总体审计};
  \node[anchor=north west,font=\figurefont\bfseries] at (88,56) {(b) 重点排查};
  \node[note,anchor=west] at (3,17) {重点类别};
  \node[note,anchor=west] at (89,17) {重点类别};
  \node[note,anchor=west] at (3,49) {常见类别};
  \node[note,anchor=west] at (89,49) {常见类别};

  \foreach \shift in {0,86}{
    \foreach \x/\y in {14/42,25/42,36/42,47/42,58/42,69/42,14/32,25/32,36/32,47/32,58/32,69/32,14/24,25/24,36/24,47/24,58/24,69/24}{
      \node[common item] at ({\shift+\x},\y) {};
    }
    \foreach \x/\y in {17/12,31/12,45/12,59/12,73/12,24/5,38/5,52/5,66/5}{
      \node[key item] at ({\shift+\x},\y) {};
    }
    \foreach \x/\y in {25/42,58/32,36/24,31/12,66/5}{
      \node[problem] at ({\shift+\x},\y) {$\times$};
    }
  }

  \foreach \x/\y in {14/42,47/42,25/32,58/32,36/24,69/24,17/12,45/12,66/5}{
    \node[audit ring] at (\x,\y) {};
  }
  \node[note,text=BookBlue,anchor=east] at (77,17) {层内随机：$\pi_h=n_h/N_h$};

  \foreach \x/\y in {25/42,47/42,58/32,69/32,36/24,58/24,31/12,45/12,66/5}{
    \node[focus ring] at ({86+\x},\y) {};
  }
  \node[note,text=BookAmber,anchor=east] at (163,17) {优先选择：风险分数$s_i\geq\tau$};

  \node[common item] at (7,-6) {};
  \node[note,anchor=west] at (11,-6) {常见类别};
  \node[key item] at (47,-6) {};
  \node[note,anchor=west] at (51,-6) {重点类别};
  \node[audit ring] at (94,-6) {};
  \node[note,anchor=west] at (99,-6) {审计入样};
  \node[problem] at (136,-6) {$\times$};
  \node[note,anchor=west] at (140,-6) {复核后确认错误};
\end{tikzpicture}

  \caption{总体审计与重点排查的抽样差异。两个面板复用同一教学总体，形状表示分层，叉号表示实际存在但抽样前未知的错误。左侧在每层随机入样并记录$\pi_h=n_h/N_h$，可按设计估计总体；右侧按冲突分数超过阈值优先选择，能提高排查命中，却改变了样本组成。图中点数仅用于说明选择机制，不按真实总体规模绘制。}
  \label{fig:knowledge-quality-assessment-correction-audit-sampling}
\end{figure*}

一个可复算的教学例子说明权重为何重要。设独立验收总体含常见类别$N_1=800$条、茶梅易混淆区域$N_2=200$条，不含开发对象$x_2$及其近重复记录。用于比较的K17和K18配置均在本次抽样前冻结，这里先看K17。两层独立进行简单随机无放回抽样，分别审计$n_1=80$条和$n_2=40$条，全部可裁决，确认K17的错误数为$4$和$8$。两层错误率为$\widehat p_1=0.05$和$\widehat p_2=0.20$，按总体占比$W_1=0.8,W_2=0.2$汇总得到$\widehat p_{\mathrm{st}}=0.08$。直接计算样本中的$12/120=0.10$会把过量抽取的重点区域看得过重；两层简单平均得到$0.125$，回答的则是“一个等权分层的平均质量”。这些数字为教学构造，后文继续用同一独立样本计算版本差异。

抽样单位必须与数据依赖一致。同一文档内的句子若共享来源和上下文，逐句简单随机抽样仍可对有限记录总体作设计推断；若流程先抽文档再审计文档内句子，文档才是第一阶段抽样单位，方差与重抽样必须保留这个簇结构。近重复照片、同一用户记录或同一生成批次也不能只因记录数很多就被当成同样多份独立证据。

样本量应由希望控制的误差范围、需要单独报告的分组、预期问题率和单例复核成本共同决定。只按总体固定比例抽取，可能在小总体上浪费预算，也可能让少数类别一个样本都没有。无法完成复核的对象同样要保留：若难例更容易无响应，删除它们会让估计偏向易判区域。重点排查可以与概率审计并行，但两批样本要有用途标识；除非为重点样本也定义了覆盖总体的入样概率和相应估计量，否则不能直接合并比例。

\section{知识质量的测量与估计}
\label{sec:knowledge-quality-assessment-correction-measurement}

审计为每个对象产生状态，测量再把状态汇总成样本级、分组级和总体级结论。这里必须区分\emph{直接核验}与\emph{有限样本估计}：复核者可以直接确认某条记录在当前参考下有错，却只能通过抽样设计估计未审计总体的错误率。指标的分母、知识版本和无法裁决比例应与估计值同时报告。

\subsection{样本级质量的测量}
\label{sec:knowledge-quality-assessment-correction-item-measurement}

单条审计的输入包括待测知识、原始对象、任务规范和参考证据，输出包括处理阶段、质量结论、问题类型、严重程度和证据。单一物种答案通常判为正确或错误；“部分正确”适用于包含多个待测属性的结果，结构化记录还要把对象、边界、类别与关系拆开。一个花卉框的位置正确而物种错误时，应分别保存这两个结论。计算二值错误率前，必须固定评价单位与计分规则，例如“对象存在、边界和类别全部合格才记正确”；不能事后把部分正确任意折算为半个错误。缺失、不适用与无法裁决按预定分母另报，不强填为正确或错误。

概率型答案要用与分布语义一致的评分。若类别集合为$Y=\{1,\ldots,K\}$，对象$i$给出概率$\bm{q}_i=(q_{i,c})_{c\in Y}$，参考类别为$y_i$，对数损失$-\log q_{i,y_i}$和Brier分数$\sum_{c\in Y}(q_{i,c}-\mathbf{1}\{y_i=c\})^2$衡量预测分布给实际结果分配了多少支持。严格适当评分规则使如实报告分布在期望意义下最优，但单个样本的分数仍不能验证概率校准\scite{gneiting2007scoring}

模型置信度、高损失和来源分歧不属于样本质量的直接裁决。它们可以触发复核，却没有独立回答原知识是否正确。类似地，文件字段齐全只能记为格式合规，来源一致只能记为意见关系。每个状态都应对应明确证据，无法裁决的记录不能悄悄从分母删除。

对构造型知识，样本级测量只报告可观察条件。正对应关系可以记录对象标识是否一致、变换是否越界；内容恢复目标可以记录目标是否来自被遮蔽部分；表示目标可以记录教师版本是否存在、输出是否退化为常数。若没有逐向量参考，就不编造“表示正确”标签，而把下游效用留给独立用途评价。

\subsection{数据集与分组质量的估计}
\label{sec:knowledge-quality-assessment-correction-dataset-estimation}

汇总前先固定计数对象。令$a_i$表示对象$i$已经获得候选知识，$j_i$表示它在审计中可以裁决，$e_i$表示可裁决知识与参考不一致，$m_i$表示目标知识缺失。条件错误率的分母是已获得且可裁决的知识，缺失率的分母是目标总体，覆盖率则要先说明“获得任意结果”还是“获得可用结果”。不同分母回答不同问题，不能只用一个百分比概括。

表~\ref{tab:knowledge-quality-assessment-correction-metrics}按分子、分母和汇总范围列出五个常用指标。类别错误、关系错配和构造违例描述不同对象，宜并列报告，不应相加成一个含义不明的总分。

\begin{table*}[htbp]
  \centering
  \small
  \begingroup
  \setlength{\tabcolsep}{2.5pt}
  \begin{tabularx}{\linewidth}{@{}>{\RaggedRight\setlength{\parindent}{0pt}\arraybackslash}p{29mm}>{\RaggedRight\setlength{\parindent}{0pt}\arraybackslash}p{43mm}>{\RaggedRight\setlength{\parindent}{0pt}\arraybackslash}p{44mm}>{\RaggedRight\setlength{\parindent}{0pt}\arraybackslash}X@{}}
    \toprule
    \textbf{指标} & \textbf{分子} & \textbf{分母} & \textbf{汇总时另报} \\
    \midrule
    覆盖条件下错误率 & 已获得且经参考确认错误的知识 & 已获得且可裁决的知识 & 无法裁决比例、参考标准与版本 \\
    总体缺失率 & 任务要求但未取得知识的对象 & 总体中应获得知识的对象 & 完全遗漏区域与不适用对象 \\
    获取覆盖率 & 达到所选接收条件的对象 & 目标总体中的适用对象 & 任意输出、可裁决或可用的接收口径 \\
    关系错配率 & 对象、方向、时间或条件错误的关系 & 已声明且可核验的关系 & 关系单位与共享对象的依赖 \\
    构造违例率 & 违反已声明生成条件的目标或关系 & 具备检查条件的已生成对象 & 变换、教师、批次与构造版本 \\
    \bottomrule
  \end{tabularx}
  \endgroup
  \caption{质量指标的分子、分母与汇总范围。未知、不适用和无法裁决状态应单列；表中指标可以并列报告，但类别错误、关系错配和构造违例描述不同对象，不能相加成含义不明的总分。}
  \label{tab:knowledge-quality-assessment-correction-metrics}
\end{table*}

当抽中的对象都能裁决时，简单随机样本的错误率可以用样本均值估计。对分层样本，设第$h$层总体量和样本量为$N_h,n_h$，总体占比$W_h=N_h/N$，层内错误率为$\widehat p_h$，则总体记录错误率估计为
\begin{equation}
  \widehat p_{\mathrm{st}}
  =
  \sum_{h=1}^{H}W_h\widehat p_h,
  \qquad
  \widehat p_h
  =
  \frac{1}{n_h}\sum_{i\in S_h}e_i.
  \label{eq:knowledge-quality-assessment-correction-stratified-rate}
\end{equation}
这个权重恢复的是记录在目标总体中的占比。若研究问题是“各类别质量是否都过关”，则应分别报告$\widehat p_h$，或明确使用按类平均；不能把两种汇总混作一个指标。

对一般不等概率抽样，已知总体规模$N$时，可以用入样概率构造Horvitz--Thompson错误率估计
\begin{equation}
  \widehat p_{\mathrm{HT}}
  =
  \frac{1}{N}
  \sum_{i\in S}\frac{e_i}{\pi_i}.
  \label{eq:knowledge-quality-assessment-correction-ht-rate}
\end{equation}
它要求目标总体内相关对象具有正入样概率，并准确保存设计概率；抽样框完全没有覆盖的区域不能靠增大权重补回。若分母也需要估计，可以采用相应的比率估计，但要说明估计目标和近似性质。Horvitz--Thompson构造及复杂抽样中的设计估计有系统论述\scite{knowledge-quality-assessment-correction-horvitz1952}

来源一致率只描述两个来源给出相同结果的频率。没有独立参考时，两个来源可以一致地错；存在多解时，它们也可能给出不同但合理的答案。来源比较因此要对齐对象、任务定义和适用范围，并把一致、冲突与无法比较分开。质量估计还要报告审计样本数、有效裁决数、无法裁决数和知识版本，使读者能够知道比例实际依赖多少证据。

\subsection{质量估计不确定性的分析}
\label{sec:knowledge-quality-assessment-correction-uncertainty}

点估计没有说明重复采用同一抽样设计时会波动多大。对从$N$条记录中简单随机无放回抽取$n$条、错误指示为$e_i$的情形，令$f=n/N$，样本方差为$s_e^2=(n-1)^{-1}\sum_{i\in S}(e_i-\widehat p)^2$，设计方差估计为
\begin{equation}
  \widehat{\operatorname{Var}}(\widehat p)
  =
  (1-f)\frac{s_e^2}{n}.
  \label{eq:knowledge-quality-assessment-correction-srs-variance}
\end{equation}
因子$1-f$是有限总体修正：审计占总体比例不可忽略时，无放回抽样比把记录视为无限总体独立抽取具有更小的抽样波动。常用有限总体、分层与整群抽样公式可参见Cochran的系统整理\scite{knowledge-quality-assessment-correction-cochran1977}

分层抽样要在层内估计波动，再按总体层权重组合。设$f_h=n_h/N_h$，层内错误指示样本方差为$s_h^2$，则
\begin{equation}
  \widehat{\operatorname{Var}}(\widehat p_{\mathrm{st}})
  =
  \sum_{h=1}^{H}
  W_h^2(1-f_h)\frac{s_h^2}{n_h}.
  \label{eq:knowledge-quality-assessment-correction-stratified-variance}
\end{equation}
式~\eqref{eq:knowledge-quality-assessment-correction-stratified-rate}和式~\eqref{eq:knowledge-quality-assessment-correction-stratified-variance}必须使用同一分层定义与抽样单位。若先抽文档再查记录，应在文档层估计第一阶段波动，或使用保留层和主抽样单位的复制权重；把所有记录放入普通自助抽样会破坏原设计。复杂调查的设计型重抽样方法也要求复制过程遵循原抽样结构\scite{knowledge-quality-assessment-correction-rao1988}

代入前面的教学数据，二值错误指示的样本方差为$s_h^2=n_h\widehat p_h(1-\widehat p_h)/(n_h-1)$，故
\[
  \begin{aligned}
    \widehat{\operatorname{Var}}(\widehat p_{\mathrm{st}})
    &=
    0.8^2(1-0.1)\frac{0.05(1-0.05)}{79}\\
    &\quad+0.2^2(1-0.2)\frac{0.20(1-0.20)}{39},\\
    \widehat{\operatorname{se}}(\widehat p_{\mathrm{st}})
    &\approx 0.02185.
  \end{aligned}
\]
若作正态近似，$8\%\pm1.96\times2.185\%$得到约$[3.72\%,12.28\%]$。它展示点估计周围的抽样波动；本例层内错误数不多，这个区间只作为近似演算，不能当作精确覆盖保证。靠近边界或用于严格验收时，应采用与设计相符的离散区间或预先规划更充分的审计。

不等概率设计还需要二阶入样概率。令$\pi_{i,j}=\mathbb{P}(i,j\in S)$且$\pi_{i,i}=\pi_i$，式~\eqref{eq:knowledge-quality-assessment-correction-ht-rate}的一个无偏设计方差估计为
\begin{equation}
  \widehat{\operatorname{Var}}(\widehat p_{\mathrm{HT}})
  =
  \frac{1}{N^2}
  \sum_{i\in S}\sum_{j\in S}
  \frac{\pi_{i,j}-\pi_i\pi_j}{\pi_{i,j}}
  \frac{e_i}{\pi_i}\frac{e_j}{\pi_j},
  \label{eq:knowledge-quality-assessment-correction-ht-variance}
\end{equation}
前提是所需的$\pi_{i,j}$为正且可得。若实际设计没有保存这些量，就应使用为该设计构造的线性化或复制权重方差，而不是退回普通独立样本公式。不等概率抽样、域估计和非响应处理的设计型框架可参见Särndal等人的论述\scite{knowledge-quality-assessment-correction-sarndal1992}

样本足够且估计不靠近边界时，可以用$\widehat p\pm z_{1-\alpha/2}\widehat{\operatorname{se}}(\widehat p)$作近似区间；小样本或错误很少时，正态近似可能越界或退化。尤其是“审计$n$条未发现错误”不等于总体无错误。若从有限总体简单随机无放回抽样并观察到零错误，给定总体中有$M$条错误时观察零错误的概率为
\begin{equation}
  \mathbb{P}(X=0\mid M)
  =
  \frac{\binom{N-M}{n}}{\binom{N}{n}}.
  \label{eq:knowledge-quality-assessment-correction-zero-hypergeometric}
\end{equation}
在置信水平$1-\alpha$下，可以反演这个超几何概率，取满足式~\eqref{eq:knowledge-quality-assessment-correction-zero-hypergeometric}不小于$\alpha$的最大$M$作为单侧上界；有放回独立Bernoulli抽样的对应上界才是$1-\alpha^{1/n}$。后者属于精确二项区间的边界情形\scite{knowledge-quality-assessment-correction-clopper1934} 不能把它不加修改地用于分层或整群设计。

抽样区间只描述由抽样造成的波动，不自动覆盖参考标签错误、审计无响应、抽样框遗漏和规范歧义。无法裁决又带来另一种不确定性：即使查遍总体，仍不知道这些对象究竟有多少错误。为明确分母，本段把$U$固定为已有候选知识的适用记录总体，大小为$N$。若可裁决结论可信，令$j_i=1$表示可以裁决，则总体错误率的\term{识别范围}为
\[
  L=\frac{1}{N}\sum_{i:j_i=1}e_i,
  \qquad
  U_{\mathrm{err}}=L+\frac{1}{N}\sum_i(1-j_i).
\]
两个端点分别把无法裁决记录视为全对与全错。只查概率样本时，按设计权重计算的$\widehat L$与$\widehat U_{\mathrm{err}}$还带有抽样误差，只是端点估计；需要覆盖保证时，应另为两端构造与抽样设计相符的联合覆盖区间，或对两端分配误差概率。例如从1000条中简单随机抽100条，全部可裁决且未见错误，两个端点估计都为0；式~\eqref{eq:knowledge-quality-assessment-correction-zero-hypergeometric}给出的95\%单侧错误数上界仍是28条，即2.8\%，不能把$[0,0]$报告为总体置信范围。

对一部分参考结果重复独立复核，可以另外检查裁决稳定性。若某一分组只审计了几条记录，合并分组会改变问题，追加审计会增加成本，暂缓结论则保留不确定性，三种选择应按用途决定。

反复查看区间直到结果达到门槛，会使固定样本量下的覆盖解释失效；同时搜索大量类别，也会增加偶然出现极端结果的机会。审计方案应预先约定主要指标、分组和停止条件，探索性发现则标明其用途并用新证据复核。样本量规划可以从希望达到的区间半宽反推，但应使用相应设计的方差、预期组内问题率和设计效应，而不是套用一个与实际抽样无关的固定比例。

\section{知识问题的检测与定位}
\label{sec:knowledge-quality-assessment-correction-detection}

总体审计回答“目标总体质量怎样”，问题检测回答“有限复核预算先看哪些对象”。检测系统输出可疑对象、触发依据、影响范围和建议核验入口，不输出已经成立的错误事实。评价检测方法时，可以比较固定复核预算下确认了多少问题、不同分组覆盖如何，以及对未入选区域随机抽查后估计遗漏多少；不能只比较风险分数大小。

\subsection{基于模型行为的问题检测}
\label{sec:knowledge-quality-assessment-correction-model-detection}

模型预测与标签分歧、异常损失、训练中反复遗忘和表示空间离群，都会把某些样本排到前面。为了减少模型直接记住待检查标签的影响，检测分数应优先来自留出预测或交叉验证预测：每个对象的分数由没有使用该对象训练的模型产生，同一对象的近重复记录也应放在同一划分单元。模型、阈值和特征若反复根据最终审计样本调整，该样本便不再独立。

% 两条边注连续排放，按整组高度为本段预留空间。
\Needspace{95mm}
这些信号都具有多义性。错误标签会使模型难以拟合，正确但罕见的样本也可能损失很高，模型能力不足或分布变化同样会产生预测分歧。基于类别条件噪声假设的confident learning可以估计和排序潜在标签问题，但其结论依赖噪声过程与概率估计条件\scite{knowledge-quality-assessment-correction-northcutt2021} 训练动态构造的数据地图也观察到难学区域经常包含标签错误，同时还包含对模型有不同作用的模糊或困难样本\scite{knowledge-quality-assessment-correction-swayamdipta2020}

图~\ref{fig:knowledge-quality-assessment-correction-detection-signal}将同一个高损失信号分解为三条候选解释。每条虚线只表示“值得核查”，右侧不同证据检查才决定知识应修正、保留，还是只限制检测模型的解释范围。

% 独立浮动页避免宽图回浮至长引文段上方，挤掉边栏的实际可用高度。
\begin{figure*}[p]
  \centering
  % 图7-2：检测信号产生候选解释，独立证据才形成裁决。
\begin{tikzpicture}[
  x=1mm,
  y=1mm,
  font=\figurefont\footnotesize,
  text=BookInk,
  signal/.style={draw=BookAmber,fill=BookAmber!9,line width=1.1pt,rounded corners=1.2mm,text width=24mm,minimum height=19mm,align=center,inner sep=1.8mm},
  explanation/.style={draw=BookMuted,fill=BookPaper,line width=0.8pt,rounded corners=1.2mm,text width=30mm,minimum height=12mm,align=center,inner sep=1.5mm},
  check/.style={draw=FigureInput,fill=FigureInput!7,line width=0.8pt,text width=30mm,minimum height=12mm,align=center,inner sep=1.5mm},
  verdict/.style={draw=FigureOutput,fill=FigureOutput!6,line width=0.8pt,text width=23mm,minimum height=12mm,align=center,inner sep=1.4mm},
  candidate/.style={-{Stealth[length=1.9mm]},draw=BookAmber,line width=0.8pt,dashed,rounded corners=1mm},
  evidence/.style={-{Stealth[length=1.9mm]},draw=BookBlue,line width=0.9pt},
  decision/.style={-{Stealth[length=1.9mm]},draw=BookTeal,line width=1.0pt},
  edge label/.style={fill=white,inner sep=0.8pt,font=\figurefont\scriptsize,text=BookMuted,align=center}
]
  \node[signal] (signal) at (15,20) {高损失或\\预测分歧\\{\scriptsize 检测线索}};

  \node[explanation] (wrong) at (58,40) {候选解释一\\知识确有错误};
  \node[explanation] (rare) at (58,20) {候选解释二\\样本正确但罕见};
  \node[explanation] (model) at (58,0) {候选解释三\\模型能力或分布不匹配};

  \node[check] (wrong-check) at (108,40) {核对原始对象、规范\\与独立参考证据};
  \node[check] (rare-check) at (108,20) {检查类别频率、覆盖\\与专业复核结果};
  \node[check] (model-check) at (108,0) {更换模型或输入条件\\比较留出行为};

  \node[verdict,draw=BookAmber,fill=BookAmber!9] (repair) at (151,40) {确认错误\\进入修正};
  \node[verdict] (keep) at (151,20) {知识可保留\\维护稀有覆盖};
  \node[verdict,draw=BookMuted,fill=white] (limit) at (151,0) {不能据此改标\\限制检测解释};

  \draw[candidate] (signal.east) -- (wrong.west);
  \draw[candidate] (signal.east) -- (rare.west);
  \draw[candidate] (signal.east) -- (model.west);
  \node[edge label,anchor=south] at (34,22) {待核查解释};

  \draw[evidence] (wrong) -- (wrong-check);
  \draw[evidence] (rare) -- (rare-check);
  \draw[evidence] (model) -- (model-check);
  \draw[decision] (wrong-check) -- (repair);
  \draw[decision] (rare-check) -- (keep);
  \draw[decision] (model-check) -- (limit);

  \node[anchor=north,font=\figurefont\scriptsize,text=BookMuted,align=center] at (84.5,-11) {同一个排序分数不能区分三种原因；裁决来自右侧检查，而不是来自分数本身。};
\end{tikzpicture}

  \caption{同一检测信号可能对应不同原因。高损失或预测分歧只把样本送入候选解释；错误知识、正确但罕见的样本和模型能力不足分别需要原始证据核验、分组与覆盖检查、模型对照。只有独立证据确认知识错误后才进入修正，箭头不表示三个解释已经由分数证实。}
  \label{fig:knowledge-quality-assessment-correction-detection-signal}
\end{figure*}
\FloatBarrier

因此，模型行为适合做排序器。输入包括原始知识、留出预测、损失和训练记录，输出是待复核队列及触发理由。错误类别、罕见类别和模型失效区域要共同进入教学或开发检查；若只用人为翻转标签验证检测器，会高估它面对真实歧义和分布变化时的辨别能力。对关系和表示目标，也可以观察配对相似度异常、输出退化或扰动响应异常，但最终仍要回到对象记录和构造条件。

阈值应与复核预算和风险相连，而不是由分数尺度自行决定。对已入选对象报告确认率，可以评估排查效率；对未入选区域保留小规模概率抽查，才能估计检测器漏掉的问题。若检测只覆盖模型可处理的输入，无法解析或模型未见过的区域还需要另一条审计路径。

\subsection{基于多源关系的问题检测}
\label{sec:knowledge-quality-assessment-correction-multisource-detection}

多来源冲突、重复对象标签不同和结构约束违例，都能形成待复核线索。检查前必须先对齐对象、语义、粒度和上下文。两条记录分别描述同一植株在春季和秋季的开花状态，它们可以不同而不矛盾；一条给物种、另一条给属，也可能只是粒度不同。只有在比较对象和成立条件相同以后，差异才构成真正冲突。

来源数量也不是证据独立性的保证。十个网页可能复制同一份错误名录，多数票只是把同一错误计算了十次；少数来源若持有原始采集记录，反而可能提供更强证据。冲突表应同时记录来源依赖、版本、时间和证据路径，不能只保留赞同票数。所有来源一致时也不能取消随机审计，因为共享错误不会产生冲突信号。

关系知识还要检查方向和条件。例如$(A,\text{采集自},B)$与$(B,\text{采集自},A)$不是同一事实，某一时间成立的属性也不能无条件复制到其他版本。结构违例可以帮助定位冲突位置，却只说明候选组合不相容；究竟哪一条有错，仍需独立材料裁决。来源很多时不必穷举所有成对比较，可以先按规范键和对象版本聚合，再优先检查影响高、来源独立性低或无法通过语义对齐解释的冲突。

\subsection{基于来源记录的问题定位}
\label{sec:knowledge-quality-assessment-correction-provenance-localization}

确认一批错误后，还要判断它们是否来自共同环节。\term{来源记录}在这里包括样本版本、标注者或模型版本、规则和词典版本、输入配置、执行状态、处理批次与时间。先从异常样本寻找共同记录，再与未受影响样本对照，随后复现候选原因并确定影响范围。错误最早被观察到的步骤不一定是产生错误的步骤：解析器可能暴露异常，而真正原因是更早的对象标识错配。

设某次词典更新后“茶梅”被大量映射为“山茶”。时间上的同时出现只能提出候选原因，因为同期样本可能换了地区，评价规范也可能改变。应在同一批冻结样本上分别运行新旧词典，保持模型和规范不变；若差异只随词典版本出现，再用原始名录核验映射。共享词典被多个规则和模型输入使用时，影响范围要扩展到全部依赖结果，而不只检查最先报警的来源。

日志缺失会直接限制结论。若无法恢复调用输入、模型版本或随机配置，只能记录“与该批次相关”的候选原因，不能宣称已经找到根因。可重现调用也不自动证明因果；仍需对照变更、独立证据或可执行检查。定位结果应包含已确认影响、可能影响和尚未覆盖的范围，为下一节选择修正粒度。

\section{知识的修正与验证}
\label{sec:knowledge-quality-assessment-correction-repair}

修正不是把检测分数最高的标签替换掉，而是根据已确认问题、可用证据和影响范围选择行动。一次处理可以改正内容、补齐遗漏、重建关系、修改构造条件，也可以在证据不足时保留不确定性。每次变更都保存原始值、候选值、采用依据、未决问题和回退入口；改变了知识或生成程序，只说明系统发生了变化，不说明质量已经改善。

\subsection{修正方案的选择与实施}
\label{sec:knowledge-quality-assessment-correction-repair-selection}

方案选择先判断三件事：当前是否有足以支持替代结果的证据，修正是否需要新增专业判断，任务定义或生成条件是否发生变化。单条操作错误可以局部修正，共享规则错误可能要求重做整批结果，任务定义变化则要建立新版本而不能与旧标准直接混算。下面四类处理覆盖“人工重判、自动执行、补充或调整、暂不强行修正”四种证据状态。

\subsubsection{基于人工复核的修正}
\label{sec:knowledge-quality-assessment-correction-human-repair}

人工修正的输入应包括原始对象、适用规范和独立证据，而不仅是一条模型建议。复核者先独立给出新判断，再查看原标签与冲突信息，输出“保留、修改、维持多解或无法裁决”及理由。操作错误可以由受训复核者处理；需要领域事实或规范解释时，应交给相应专家，并把可推广的裁决写回规范。

对象错配说明单条修正可能牵动多条记录。若照片S042被写到了植株P017名下，复核者要先恢复照片的真实对象，再检查S042的类别、与同株照片的正对应关系以及以P017为锚点产生的邻居结果。只把类别改对会留下错误关系。修正范围应沿来源依赖展开，但每条受影响记录仍要保存自己的变更依据。

独立复标与知情裁决服务于不同目的。前者减少已有答案的锚定，适合产生新的判断；后者把各方意见和证据放在一起，适合解释争议。多人仍然分歧时，可以保留答案集合和条件，不必无限追加人数追求单一标签。何时停止取决于使用风险、信息增益和成本，不能用“已经复核一次”代替这些条件。

\subsubsection{基于自动推断的修正}
\label{sec:knowledge-quality-assessment-correction-automatic-repair}

自动修正只有在转换规则、外部证据或经过验证的预测能够支持替代结果时才成立。类别编码从旧编号无歧义地映射到新编号，属于限定条件下保持语义的确定性转换；模型对疑似错标样本重新预测，则只是生成候选。两者不能共用“置信度高于阈值即修正成功”的标准。

一个可审计流程依次生成候选、检查适用条件、决定自动接收或转人工、保存前后差异。确定性规则要说明为何保持对象和语义；预测替换要分别记录模型分数、替换风险和复核要求。模型与原标签若来自相同训练数据或共享来源，多次预测一致仍可能重复同一错误。由同一教师重新生成表示目标，也不构成对教师输出的独立证明。

自动方案应先在开发样本上确定规则与阈值，冻结后再进入留出审计。小范围试用要同时检查被修改对象和未修改对象：前者暴露改错风险，后者暴露规则漏修与作用边界。出现系统性副作用时回退到旧版本，而不是只继续提高阈值直到开发样本看起来正确。

\subsubsection{基于问题类型的补充与调整}
\label{sec:knowledge-quality-assessment-correction-type-repair}

不同问题要求改变不同对象。遗漏需要补齐知识，错配需要重建对象对应，粒度不足需要增加局部信息，定义歧义需要修订规范并重查受影响记录，视图破坏语义则需要改变变换和目标构造。前四类可能仍保留原任务定义；定义变化会改变“正确”的含义，必须另建版本并说明新旧口径的可比部分。

表~\ref{tab:knowledge-quality-assessment-correction-repair-scope}从已确认问题出发确定修正范围。使用它的前提是问题已经由参考证据核实；若输入仍只是检测分数或来源冲突，应先返回问题定位，而不能直接选择批量替换。

\begin{table*}[htbp]
  \centering
  \small
  \begingroup
  \setlength{\tabcolsep}{2.5pt}
  \begin{tabularx}{\linewidth}{@{}>{\RaggedRight\setlength{\parindent}{0pt}\arraybackslash}p{26mm}>{\RaggedRight\setlength{\parindent}{0pt}\arraybackslash}p{47mm}>{\RaggedRight\setlength{\parindent}{0pt}\arraybackslash}p{43mm}>{\RaggedRight\setlength{\parindent}{0pt}\arraybackslash}X@{}}
    \toprule
    \textbf{核查结果} & \textbf{改变对象与范围} & \textbf{可选处理} & \textbf{后续核验} \\
    \midrule
    单点错标 & 当前答案及引用它的派生记录 & 独立复标后替换或保留多解 & 复核改后结果，抽查关联记录 \\
    对象错配 & 对象绑定、同对象标签与关系 & 恢复对象后重建关联结果 & 对象、方向与影响范围 \\
    信息缺失 & 缺失知识及同类、同批次对象 & 补标、追加获取或保留弃权 & 覆盖变化与新增部分质量 \\
    定义变化 & 规范版本及旧定义下的相关记录 & 分版本保留、有限映射或重标 & 共同定义域内的前后差异 \\
    构造条件失效 & 来源规则及同配置生成的目标 & 修改规则，重生成受影响批次 & 构造条件与独立下游用途 \\
    证据不足 & 状态、候选集合与使用方式 & 暂缓、降权、隔离或追加证据 & 覆盖损失与再次触发条件 \\
    \bottomrule
  \end{tabularx}
  \endgroup
  \caption{核查结果与处理范围的对应。前五行处理已有证据支持的问题，末行处理仍无法裁决的情况；高损失或来源冲突本身不能确立问题类型。}
  \label{tab:knowledge-quality-assessment-correction-repair-scope}
\end{table*}

例如，裁剪规则删除了花朵主体，却仍把裁剪图与原图记为正对应。修正需要限制裁剪区域、重新生成关系，并检查使用旧关系训练或评价的派生版本。补标和重生成会改变覆盖与数据分布，因此审计总体也要更新；只看已改样本会漏掉新增覆盖是否偏向某些类别。

\subsubsection{无法可靠修正时的处理}
\label{sec:knowledge-quality-assessment-correction-unresolved}

证据不足时，诚实保留不确定性比制造确定答案更可靠。可选行动包括保留多个合理标签或关系、标记待复核、降低训练权重、暂时隔离，以及等待新证据。它们改变的是知识的表达或参与训练的方式，不一定改变内容本身。降低权重不会让错误变正确，删除记录也只是移除了风险及相应覆盖。

少数类别持续得到低置信分数时，批量删除尤其危险。低分可能来自模型缺少该类别能力；删除后，训练数据会进一步失去这一任务区域。应先抽查独立证据，再比较保留候选集合、允许弃权和追加专业复核的代价。若当前任务不能容忍未决状态，可以暂时隔离，但要报告覆盖下降，不能把剩余数据的较高正确率解释成原总体质量提高。

未决记录需要有去向：保存原始内容和候选证据，规定来源更新、达到风险阈值或积累到一定数量时重新触发审计，并定期检查它们是否集中在特定群体。停止追加证据也应有明确理由，例如使用风险低于门槛或边际复核成本过高，而不是让“待复核”永久成为不可见状态。

\subsection{修正效果的验证}
\label{sec:knowledge-quality-assessment-correction-repair-validation}

验证包含两个不同问题。修正结果的独立复核问“这些已经改动的对象现在是否正确”；修正策略的留出验证问“冻结后的处理规则遇到未参与设计的新对象时是否仍有效”。前者可以复核同一对象，但判断证据或复核者要与修改过程适当隔离；后者必须把策略开发样本与评价样本隔离，不能在查看留出结果后继续调规则还沿用原有独立性声明。

图~\ref{fig:knowledge-quality-assessment-correction-evidence-isolation}把两类隔离画成上下两条路径。上方允许同一对象形成修正前后配对，隔离的是判断证据；下方隔离的是策略开发与推广评价，冻结边界之后不能再把留出结果送回策略选择。

\subsubsection{修正结果的独立验证}
\label{sec:knowledge-quality-assessment-correction-independent-validation}

每个对象应保留修正前后成对记录。设同一对象$i$在修正前后的错误指示为$e_i^{(0)}$和$e_i^{(1)}$，定义改进量$d_i=e_i^{(0)}-e_i^{(1)}\in\{-1,0,1\}$。$d_i=1$表示错误被纠正，$d_i=-1$表示原本正确却被改错，$d_i=0$还包含“始终正确”和“始终错误”两种情况，报告时应把四种转移分别列出。

要把配对变化推广到总体，比较的两个版本或对总体统一执行的修正策略，必须在评价抽样之外确定。评价样本不能用于选择哪些对象应修改，也不能用于调整修正规则；策略可以按对象特征区别处理，但判定规则须预先冻结。对这样固定的前后结果，在各层独立进行简单随机无放回审计时，总体错误率下降量的估计及其方差为
\begin{equation}
  \widehat\Delta_{\mathrm{st}}
  =
  \sum_{h=1}^{H}W_h\overline d_h,
  \qquad
  \widehat{\operatorname{Var}}(\widehat\Delta_{\mathrm{st}})
  =
  \sum_{h=1}^{H}W_h^2(1-f_h)\frac{s_{d,h}^2}{n_h},
  \label{eq:knowledge-quality-assessment-correction-paired-difference}
\end{equation}
其中$\overline d_h$和$s_{d,h}^2$是第$h$层配对差值的均值与样本方差。这里直接分析同一对象的差值，保留了前后相关性；把两次错误率当作独立样本会使用错误的方差。简单随机二值配对中，若$b$条由错变对、$c$条由对变错，则$\widehat\Delta=(b-c)/n$；McNemar检验利用这两类不一致对检验边际比例是否相同，但效应大小和区间仍应报告\scite{knowledge-quality-assessment-correction-mcnemar1947}

随机入样本身并不足够。假设1000条原记录全部错误，随机抽100条后只改正这100条，再请独立专家确认全部改对。若把样本中的$d_i=1$直接代入上式，会误报总体下降100个百分点、方差为0；实际只下降10个百分点。这里修改对象由入样决定，评价样本已参与修正。应报告这100条的确认纠正数及其占总体的已知份额；若要估计修改后总体或策略推广效果，应另取独立评价样本。

设计匹配仍然适用。按簇抽样时，配对差值要在簇层估计或重抽样；只评价高风险范围时，结论也限于该范围。若策略预先规定只修高风险对象，而评价样本覆盖整个总体，则可以估计全体的变化，包括未处理区域的零变化。规范改变时，只有新旧定义共同覆盖的对象可以作同口径配对；其余对象应分版本报告。

\begin{figure*}[htbp]
  \centering
  % 图7-3：结果独立复核与策略留出验证使用不同的隔离要求。
\begin{tikzpicture}[
  x=1mm,
  y=1mm,
  font=\figurefont\footnotesize,
  text=BookInk,
  record/.style={draw=FigureInput,fill=FigureInput!8,line width=0.8pt,text width=21mm,minimum height=13mm,align=center,inner sep=1.4mm},
  action/.style={draw=FigureModel,fill=FigureModel!8,line width=0.9pt,rounded corners=1.2mm,text width=23mm,minimum height=13mm,align=center,inner sep=1.4mm},
  evidence/.style={draw=BookAmber,fill=BookAmber!8,line width=0.8pt,text width=25mm,minimum height=12mm,align=center,inner sep=1.4mm},
  result/.style={draw=FigureOutput,fill=FigureOutput!6,line width=0.9pt,text width=22mm,minimum height=13mm,align=center,inner sep=1.4mm},
  flow/.style={knowledge arrow,draw=BookTeal,line width=1.0pt,rounded corners=1mm},
  evidence flow/.style={knowledge arrow,draw=BookAmber,line width=0.9pt,rounded corners=1mm},
  boundary/.style={draw=BookMuted,line width=0.8pt,dashed},
  note/.style={font=\figurefont\scriptsize,text=BookMuted,align=center}
]
  \draw[BookRule,line width=0.7pt] (0,37) -- (166,37);
  \node[anchor=north west,font=\figurefont\bfseries] at (1,93) {(a) 已改结果的独立复核：对象可以相同，判断证据需要独立};
  \node[record] (before) at (14,58) {开发对象$x_2$\\修正前记录};
  \node[action] (repair) at (48,58) {依据开发证据\\实施修正};
  \node[record,draw=FigureModel] (after) at (82,58) {开发对象$x_2$\\修正后记录};
  \node[action,draw=BookAmber,fill=BookAmber!8] (review) at (122,58) {未参与修改者\\独立复核};
  \node[result] (verdict) at (153,58) {确认改对／改错\\或无法裁决};
  \node[evidence] (new-evidence) at (122,79) {新证据或盲化材料};

  \draw[flow] (before) -- (repair);
  \draw[flow] (repair) -- (after);
  \draw[flow] (after) -- (review);
  \draw[flow] (review) -- (verdict);
  \draw[evidence flow] (new-evidence) -- (review);

  \node[anchor=north west,font=\figurefont\bfseries] at (1,32) {(b) 总体或策略效果评价：在评价抽样前冻结策略};
  \node[record] (development) at (15,9) {开发审计样本\\已查看问题};
  \node[action] (design) at (50,9) {设计规则、模型\\或接收阈值};
  \node[action,draw=BookInk,fill=BookPaper] (freeze) at (85,9) {冻结修正策略\\和评价口径};
  \node[evidence] (holdout) at (123,9) {独立抽样审计\\比较前后版本};
  \node[result] (effect) at (153,9) {策略效果\\及适用范围};

  \draw[flow] (development) -- (design);
  \draw[flow] (design) -- (freeze);
  \draw[flow] (freeze) -- (holdout);
  \draw[flow] (holdout) -- (effect);
  \draw[boundary] (104,-4) -- (104,22);
  \node[note,anchor=south] at (104,22) {评价边界};
  \node[note,anchor=north,text width=162mm] at (83,-6) {查看评价结果后若继续调整策略，需另取独立证据；开发对象不进入本例的验收总体。};
\end{tikzpicture}

  \caption{修正与验证的证据隔离。上面板复核同一被修改对象，独立性来自未参与修改的复核者或新证据；下面板用开发审计确定策略，在评价抽样前冻结，再用留出审计评价。上面板只确认已改结果；用于总体推断的配对样本还须满足下面板的隔离条件。}
  \label{fig:knowledge-quality-assessment-correction-evidence-isolation}
\end{figure*}

独立复核要统计纠正、改错、未解决和退出使用四类结果，并抽查未修改区域。只复核被改对象会看见修正精度，却看不见规则漏掉的问题；只看总体平均又可能掩盖少数类别恶化。关系重建要分别检查对象和条件，表示目标更新则要分别检查构造条件与冻结下游用途。证据不足或关键分组下降时，应回退、增加审计或限制使用范围。

\subsubsection{修正收益与成本的比较}
\label{sec:knowledge-quality-assessment-correction-benefit-cost}

“修改正确”不等于“值得继续修改”。方案比较要同时看纠错、新增错误、覆盖、关键分组、下游用途和完整投入。固定数据版本、训练配置和评价集以后，才可以把差异归到修正内容；若同时更换模型、增加训练预算和扩大数据，不能只把收益归因于标签修正。

现在把K17的开发诊断接到一次具体验收。开发复核在$x_2$等对象上确认共享旧词典失效，因此K18暂时把这两条规则的权重置零，其他作答来源仍按等权合并；无可用答案或平票时保留未决。这是一种隔离失效来源的处置，尚未修好词典本身。在$x_2$上，模型的“茶梅”成为剩下的有效票。两版均保留原始来源输出，便于以后恢复和比较。

策略、分层及门槛在打开前述1000条独立验收总体的参考答案之前固定，并在评价抽样前冻结。K18按同一规则重算全部1000条，不只重算抽中的120条。教学记录显示两版都给出1000条可接收的类别答案，没有新增、退出或未决对象；这里的接收指满足输出协议与阈值，不表示已经确认正确。独立复核同一组80条和40条样本，得到表~\ref{tab:knowledge-quality-assessment-correction-benefit-cost}中的配对转移。

常见层有3条改对、1条改错，易混淆层有6条改对、2条改错，因此
\[
  \overline d_1=\frac{3-1}{80}=0.025,\qquad
  \overline d_2=\frac{6-2}{40}=0.10,
\]
\[
  \widehat\Delta_{\mathrm{st}}=0.8(0.025)+0.2(0.10)=0.04.
\]
加权纠正率为6\%，新增错误率为2\%，净改善为4个百分点。K18的分层错误数为2和4，总体错误率点估计为$0.8(2/80)+0.2(4/40)=4\%$，低于K17的8\%。样本中仍有3条原错误未解决，另有3条新错误；把二者分开，才能看见策略漏修和误替换这两种代价。

配对方差不能用两次错误率方差简单相加。由$d_i^2$只在改变对错的对象上等于1，得到
\[
  s_{d,1}^2=\frac{3+1-80(0.025)^2}{79}=0.05,\qquad
  s_{d,2}^2=\frac{6+2-40(0.10)^2}{39}\approx0.19487.
\]
代入式~\eqref{eq:knowledge-quality-assessment-correction-paired-difference}得到标准误约0.02271。正态近似95\%区间为
\[
  0.04\pm1.96(0.02271)
  \approx[-0.00452,0.08452].
\]
用百分点表示，改善量约落在$[-0.45,8.45]$之间。改变对错的配对只有12条，近似误差不可忽略；这个演算不足以提供严格的95\%发布保障，也没有支持“改善确定为正”。

\begin{table*}[htbp]
  \centering
  \small
  \begingroup
  \setlength{\tabcolsep}{4pt}
  \begin{tabular*}{\linewidth}{@{\extracolsep{\fill}}lrrrrrr@{}}
    \toprule
    \textbf{分层} & \textbf{抽查数} & \textbf{错$\to$对} & \textbf{对$\to$错} & \textbf{对$\to$对} & \textbf{错$\to$错} & \textbf{退出} \\
    \midrule
    常见类别 & 80 & 3 & 1 & 75 & 1 & 0 \\
    易混淆区域 & 40 & 6 & 2 & 30 & 2 & 0 \\
    \bottomrule
  \end{tabular*}

  \medskip

  \begin{tabularx}{\linewidth}{@{}>{\RaggedRight\setlength{\parindent}{0pt}\arraybackslash}p{29mm}>{\RaggedRight\setlength{\parindent}{0pt}\arraybackslash}X>{\RaggedRight\setlength{\parindent}{0pt}\arraybackslash}p{56mm}@{}}
    \toprule
    \textbf{验收项} & \textbf{预先给定的要求} & \textbf{本次结果} \\
    \midrule
    错误改善 & 点估计至少2个百分点，且区间下界大于0 & 点估计4个百分点；近似下界为负，证据不足 \\
    分组表现 & 两层错误率点估计均不升高 & 5\%降至2.5\%；20\%降至10\% \\
    接收覆盖 & 总体及每层均不低于99\% & 两版均为100\%；退出为0 \\
    本轮投入 & 人工不超过32工时，计算不超过1机时 & 开发4、复核24、维护2工时；重算0.5机时 \\
    \bottomrule
  \end{tabularx}
  \endgroup
  \caption{K17到K18的教学验收。上表是独立样本中的原始配对数，汇总需使用总体层权重；下表连接预定要求与结果。全部审计对象均可裁决，且没有退出，故四类对错转移覆盖全部样本。点估计改善而不确定性仍大，本次暂缓发布。}
  \label{tab:knowledge-quality-assessment-correction-benefit-cost}
\end{table*}

下游收益只在知识实际用于任务时解释。类别修正可以比较固定训练流程下的评价指标，关系重建可以检查关系任务及其派生模型，表示目标则应在冻结的下游任务和明确使用方式中比较。构造违例减少是必要证据，却不能替代用途评价。人工复核、自动调用、重生成、再训练、再次审计和长期维护都属于成本；继续修正的边际收益低于新增知识或保留现状时，停止也可以是有依据的选择。

本例的覆盖与投入达到门槛，但错误改善的证据未达到要求，决定暂缓发布K18。复核投入24工时按120条、每条12分钟计算，另计开发4工时和维护2工时，共30工时；两类资源没有强行换成同一金额。此轮尚未重训下游模型，重训成本与用途收益均未发生，不能把标签错误率下降当作已取得的模型收益。若继续推进，应预先设计新的审计批次及合并规则，或重新开发策略后另作独立验收，不能不断加看样本直到区间过线。

\subsection{知识版本与持续质量的维护}
\label{sec:knowledge-quality-assessment-correction-version-maintenance}

质量结论只对特定知识、规范和来源版本成立。一个可复查版本至少包含数据范围、知识内容或目标构造配置、来源依赖、修改理由、审计设计、参考依据和验证结果。只有一个递增版本号，却无法回答它由哪些活动和来源生成，仍不足以复现判断。数据集文档应随维护过程更新，而不是只在首次发布时写一次\scite{gebru2021datasheets}

静态标签和关系要保存新增、删除、替换与未决状态的差异；动态生成的表示目标还要保存模型、输入变换、随机配置和能否重现数值。W3C PROV数据模型把实体、活动、责任主体和派生关系分开，为记录“哪个版本由什么过程和来源产生”提供了通用表达\scite{knowledge-quality-assessment-correction-w3c2013provdm} 本章不要求采用某一存储系统，但这些关联不能只存在于临时日志。

图~\ref{fig:knowledge-quality-assessment-correction-version-feedback}展示共享来源更新后的追踪路径。实线让验证记录绑定到具体知识版本，虚线只触发复查、重生成和回退；新版本通过验证以前，旧版本仍保留为可恢复基线。

\begin{figure*}[htbp]
  \centering
  % 图7-4：来源、构造、知识和验证记录之间的可追溯版本依赖。
\begin{tikzpicture}[
  x=1mm,
  y=1mm,
  font=\figurefont\footnotesize,
  text=BookInk,
  data/.style={draw=FigureInput,fill=FigureInput!8,line width=0.8pt,text width=20mm,minimum height=12mm,align=center,inner sep=1.4mm},
  source/.style={draw=FigureModel,fill=FigureModel!8,line width=0.9pt,text width=20mm,minimum height=12mm,align=center,inner sep=1.4mm},
  config/.style={draw=BookMuted,fill=BookPaper,line width=0.8pt,text width=20mm,minimum height=12mm,align=center,inner sep=1.4mm},
  knowledge/.style={draw=BookTeal,fill=BookTeal!7,line width=1.0pt,text width=23mm,minimum height=12mm,align=center,inner sep=1.4mm},
  validation/.style={draw=BookAmber,fill=BookAmber!8,line width=0.9pt,text width=24mm,minimum height=12mm,align=center,inner sep=1.4mm},
  release/.style={draw=FigureOutput,fill=FigureOutput!6,line width=0.9pt,text width=19mm,minimum height=12mm,align=center,inner sep=1.4mm},
  dependency/.style={knowledge arrow,draw=BookMuted,line width=0.9pt,rounded corners=1mm},
  action/.style={knowledge arrow,draw=BookAmber,line width=0.9pt,dashed,rounded corners=1mm},
  publish/.style={knowledge arrow,draw=BookTeal,line width=1.0pt},
  edge label/.style={fill=white,inner sep=0.8pt,font=\figurefont\scriptsize,text=BookMuted,align=center}
]
  \node[data] (raw) at (10,20) {原始数据\\版本D5};
  \node[source] (source-old) at (40,42) {原来源配置\\版本M3};
  \node[source,draw=BookAmber] (source-new) at (40,-1) {隔离失效规则\\配置M4};
  \node[config] (config) at (75,20) {对象与输入\\配置C2};

  \node[knowledge,draw=BookAmber] (knowledge-old) at (75,42) {知识版本K17\\保留旧路径};
  \node[knowledge,draw=BookAmber] (knowledge-new) at (75,-1) {知识版本K18\\重生成候选};

  \node[validation] (validation-old) at (111,42) {验证记录R17\\绑定K17};
  \node[validation] (validation-new) at (111,-1) {新审计记录R18\\绑定K18};
  \node[release] (release-old) at (148,42) {已发布K17\\可回退};
  \node[release] (release-new) at (148,-1) {候选K18\\本例暂缓发布};

  \draw[dependency] (raw.east) -- ++(4,0) |- (source-old.west);
  \draw[dependency] (raw.east) -- ++(4,0) |- (source-new.west);
  \draw[dependency] (source-old) -- (knowledge-old);
  \draw[dependency] (config.north) -- (knowledge-old.south);
  \draw[dependency] (source-new) -- (knowledge-new);
  \draw[dependency] (config.south) -- (knowledge-new.north);
  \draw[dependency] (knowledge-old) -- (validation-old);
  \draw[dependency] (knowledge-new) -- (validation-new);
  \draw[publish] (validation-old) -- (release-old);
  \draw[publish] (validation-new) -- (release-new);

  \draw[action] (source-old.east) -- ++(5,0) |- (source-new.east);
  \draw[action] (knowledge-old.east) -- ++(5,0) |- (knowledge-new.east);
  \draw[action] (release-new.east) -- ++(3,0) |- (release-old.east);
  \node[edge label,text=BookAmber,anchor=east] at (54,20.5) {来源更新};
  \node[edge label,text=BookAmber,anchor=west] at (95,20.5) {重生成};
  \node[edge label,text=BookAmber,anchor=east] at (160,20.5) {验证失败回退};
  \node[font=\figurefont\scriptsize,text=BookAmber,anchor=south,align=center] at (75,51) {来源配置改变后，保留K17，并重算受影响记录形成K18};

  \node[anchor=north,font=\figurefont\scriptsize,text=BookMuted,align=center] at (82,-13) {实线：版本生成依赖；虚线：反馈动作，具体动作标在对应边旁。};
\end{tikzpicture}

  \caption{知识版本、来源依赖与质量反馈。实线表示原始数据、来源版本和构造配置对知识版本及验证记录的生成依赖；虚线表示来源更新后触发复查、重生成或回退。旧版本与新版本均保留，验证记录绑定具体知识版本；复查反馈用于形成新证据，不能把同一审计结果无限重复当作独立验证。}
  \label{fig:knowledge-quality-assessment-correction-version-feedback}
\end{figure*}

来源模型、规则词典或任务规范更新时，维护流程先沿依赖记录确定受影响对象，再对冻结样本运行新旧配置，执行局部概率审计，比较版本差异，最后决定发布、限制范围或回退。定期抽查用于发现缓慢退化，事件触发检查用于响应来源和定义变化。连续监测仍然只产生线索；一旦监测样本参与阈值和方案选择，就要准备新的独立证据验证发布结论。

章首$x_2$的开发排查确认共享来源错误，K18因此隔离两条失效规则，并用固定配置重算总体。独立验收虽然得到4个百分点的净改善点估计，仍因证据不足暂缓发布；版本记录应保留这个决定和未解决的错误，不能只保存更好看的点估计。暂缓K18也不等于重新证明K17可靠，已确认的旧版问题仍应保留风险处置与复查安排。

下一步若选择修好词典，就形成新的来源配置和待验收版本；若发现照片与植株绑定错误，则返回知识对齐；若问题来自模型输出或裁剪条件，则返回相应获取环节。反馈返回哪个步骤，由已确认的原因决定。新的开发过程使用过的证据要继续标明用途，后续发布判断另有独立评价依据。

版本管理还要保留覆盖变化。修正删除了难例、规范缩小了适用范围或新来源补入了少数类别时，平均错误率的变化可能来自总体组成改变。发布记录应同时报告共同对象上的配对变化、新增与退出对象、目标总体定义及仍未解决的问题。这样，质量反馈才能准确返回知识获取、来源选择、语义对齐或融合环节，而不会把每次版本变化都误写成进步。

\section{本章小结}
\label{sec:knowledge-quality-assessment-correction-summary}

知识质量不能脱离标准、范围和证据来判断。任务答案需要与独立参考核验，关系知识需要检查对象及成立条件，表示目标则先核查构造条件，再在独立用途上评价。总体审计通过概率抽样和设计匹配的权重、方差回答“整体怎样”；重点排查通过风险信号提高发现问题的效率，却不能把命中比例直接推广到总体。

检测结果是复核入口，不是错误证明。模型高损失、来源冲突和异常日志都可能有多种解释，只有回到原始对象、任务规范与独立证据才能确认问题。修正以后，同一对象的前后变化应按配对设计分析；已改结果的独立复核检查修改是否正确，冻结策略的留出验证则检查方法能否推广，两者也不能互换。

最终质量决策要同时考虑错误、覆盖、重要分组、下游用途、不确定性和成本。抽样区间不能掩盖参考错误与系统遗漏，平均改善也不能抵消少数区域恶化。保存来源、配置、知识版本和验证记录，才能把已确认问题反馈到获取、对齐或融合过程，并让后来者知道哪些结论已经得到支持、哪些仍然只是一条线索。
